% Chapter 5

\chapter{Introduction} % Write in your own chapter title
\label{Chapter5}
\lhead{Chapter 5. \emph{Introduction}} % Multi-label classification


\section{Overview}

In machine learning literature, multi-label classification is a variant of the classification problem in which each instance has several labels. In general, the task of multi-label learning is to find a model that can maps inputs \textbf{x} to binary vector \textbf{y}. While it is different with multi-class classfication in terms of the later one scalars outputs in classification problem.

\graphicspath{ {./Figures/} }
\begin{figure}[!htb]
    \centering
	\includegraphics[width=0.8\textwidth]{beach.jpg}
    \caption{Example Image}%
    \label{fig:MultilableImage}%
\end{figure}

The difference between single-label classification and multi-label classification is the outputs of datasets.
In Figure \ref{fig:MultilableImage}, we can classify it as a picture of a beach in classification literature $\in \{Yes, No\}$. In multi-label literature, we can tag sea, beach, chairs, sky, cloud for the picture $\in \{beach, sea, chairs, sand\}$.

\begin{table}
\centering
\begin{tabular}{|c c c c c | c c c c|}
\hline
$X_{1}$ & $X_{2}$ &  $X_{3}$ & $X_{4}$ & $X_{5}$ & $Y_{1}$ & $Y_{2}$ & $Y_{3}$ & $Y_{4}$ \\
\hline
1 & 0.4 & 0.2 & 0 & 1 & 0 & 1 & 1 & 0 \\
0 & 0.2 & 0.6 & 1 & 1 & 1 & 0 & 1 & 0 \\
0 & 0.5 & 0.8 & 1 & 1 & 0 & 0 & 1 & 0 \\
1 & 0.1 & 0.4 & 0 & 1 & 1 & 1 & 1 & 0 \\
1 & 0.8 & 0.2 & 0 & 1 & 0 & 1 & 1 & 0 \\
\hline
0 & 0.5 & 0.3 & 1 & 1 & ? & ? & ? & ? \\
\hline
\end{tabular}
\caption{Multilabel $Y_{1},...,Y_{L} \in 2^L$}
\ref{tab:MultilabelTable}
\end{table}

In general, there exists two main approaches of tackling the multi-label classification problem. One transfers a multi-label problem into a set of binary classification problems which can be handled by a set of binary classifiers. The other one applies algorithms to complete multi-label classification. 

Several problem transformation methods can be applied to multi-label classification. A baseline method, named the binary relevance method\citep{read2011classifier}, trains one binary classifier for each label independently. Depending on the results of the classifiers, the combined model predicts all labels for an unseen sample. The method divides the problem into multiple binary works which has something in common with the one-vs-all method for multi-class classification. Problem transformation benefits from scalability and flexibility because single-label classifier can be implemented to job. Support Vector Machine, Naive Bayes, $k$ Nearest Neighbor have been used in the method\citep{read2011classifier}.

There are some other transformation methods, for instance label powerset transformation. The method builds up one binary classifier for each label combination verified in the training dataset\citep{tsoumakas2006multi}. The random $k$-labelsets algorithm\citep{tsoumakas2007random} utilizes multiple label powerset classifier which is trained on a random sub dataset of the labels. A voting scheme is used to predict unseen samples via a ensemble method.

Multilayer neural network can be trained to learn a model from multi-label examples. We will generate a 3 labels dataset and use a 2 layer network to do classification.

\section{Multi-Label Learning}
Let $X = R^d$ represent the domain of dataset and let $Y = {1,2,..,L}$ be the finite set of labels. Given that we have a set of training dataset $T = {(x_{1},Y_{1}),(x_{2},Y_{2}),...,(x_{l},Y_{l})} (x_{i} \in X, Y_{i} \subset Y)$ which are extracted from an unknown distribution $D$. Our target of task is to learn a multi-label classifier $h: X \to 2^y $ via optimizing specific evaluation metric. However, instead of learning a multi-label classifier, we will learn a function $f$ while $f(X) \to R^d$. Supposing that a high performance classifier can output a closer subset for labels in $Y_{i}$ than those missing or exceeding in $Y_{i}$, which means $f(x_{i}, y_{1}) > f(x_{i}, y_{2})$ if $y_{1} \in Y{i}$ and $y_{2} \notin Y{i}$. We can transfer real valued function $f(\cdot , \cdot)$ to a ranking function $r(\cdot , \cdot)$ that maps the outputs of $f(x_{i}, y_{1})$ to any $y \in Y$ if $f(x_{i}, y_{1}) > f(x_{i}, y_{2})$. It is worth noting that the multi-label classifier $h(\cdot)$ can be derived from the function $f(\cdot , \cdot)$ where $h(x_{i}) = {y|f(x_{i}, y) > t(x_{i}), y \in Y}$, and $t(\cdot)$ is a threshold function.

Single-label and multi-class classification can be regarded as two degenerated variants of multi-label learning problem if each sample has only one single label. However, multi-label problem is much more difficult than traditional single-label problems because of high dimensional output space. For example, the number of label sets increases exponentially with increasing number of class labels. If there are $10$ class labels for dataset, there are $2^20$ possible label sets maximun.

The challenge of huge combination of output labels is hard to overcome. One of methods is to investigate dependency among labels. For example, if an image labelled with $castle$, it would be highly labelled with $brick$ and $mountain$. A movie which is labelled with $comedy$ is unlikely related to a $documentary$. Therefore, successful exploitation of the label correlations is regarded as an effective approach to high accuracy multi-label learning process. There are three categories, based on the order of correlation of labels, for existing strategies to find the relation between labels. They are first order strategy, second order strategy and high-order strategy. 

The first order strategy treats label by label independently and ignores correlation between labels. It can be regarded as decomposing a multi-label learning problem into a set of binary classification problems based on each label. The method benefits from simple computation and high efficiency. However, accuracy could be suboptimal because of ignoring the correlations of labels.

The second order strategy considers pairwise relations between labels, for example, interaction between any pair of labels, or the ranking between related labels and unrelated labels. The method can help to achieve good generalization performance because the label correlations are investigated in some extent. In real world, there are higher order correlations than second order assumption in many applications.

The high order strategy investigates more than $2$ order correlations among labels which can be the influences on each label or addressing connections among sub space of output labels. It is obvious that high order strategy has the strongest model capabilities than previous two strategies on the cost of complexity and intensive computation.

\section{Evaluation Metrics}

In supervised learning settings, different metrics, like accuracy and area under the ROC curve, are used to evaluate the generalization performance of learning system. In multi-label learning, evaluating performance is more complicated  than single-label problems with the increasing number of labels simultaneously. Therefore, two main types of evaluation methods are implemented in multi-label learning, say example-based metrics\citep{ghamrawi2005collective} and label-metrics\citep{tsoumakas2007random}.

The two types of metrics evaluate the outputs of classifier from different perspective. Given that $S = {(x_{i}, Y_{i})}$ is the test instance and $h(\cdot)$ is the learned multi-label classifier. Example-based metrics achieve all class labels of each test example, and then compute the mean value of test set to evaluate generalization performance. Compared to considering all class labels simultaneously, label-based metrics evaluate performance by treating each class label separately and computing macro/micro-averaged value of all class labels.

In a supervised classification problem, there have a ground truth output and a predicted output of test instance. So the results of each test instance can be be assigned to one of the four categories:
\begin{itemize}
\item True Positive (TP) - label is positive and prediction is also positive
\item True Negative (TN) - label is negative and prediction is also negative
\item False Positive (FP) - label is negative but prediction is positive
\item False Negative (FN) - label is positive but prediction is negative
\end{itemize}

Here we define a set $D$ of $N$ instances and $Y_{i}$ to be a family of ground truth label sets and $P_{i} = h(x_{i})$ to be a family of predicted label set. The set of all unique labels is
\begin{equation}\label{eq:UniLabel}
L = \bigcup_{i=0}^{N-1} L_{i}
\end{equation}
While the definition of indicator function $I_{A}$ on a set $A$ is presented as
\begin{equation}\label{eq:IndicatorFunc}
I_{A}(x) =
  \begin{cases}
    1       & \quad \text{if } x \in A\\
    0  & \quad \text{otherwise}\\
  \end{cases}
\end{equation}

\subsection{Example-based Metrics}

\textbf{Hamming Loss} evaluates performance via counting the number of misclassification labels. The smaller value of hamming loss, the better performance.
\begin{equation}\label{eq:HammingLoss}
\frac{1}{N \cdot \left|L\right|} \sum_{i=0}^{N - 1} \left|L_i\right| + \left|P_i\right| - 2\left|L_i
          \cap P_i\right|
\end{equation}

\textbf{Subset Accuracy} evaluates the fraction of correctly predicted instance while the predicted label set is identical to the ground truth label set. It is equivalent to traditional accuracy metric.
\begin{equation}\label{eq:SubsetAcu}
\frac{1}{N} \sum_{i=0}^{N-1} I_{\{L_i\}}(P_i)
\end{equation}

\textbf{Precision} is defined as
\begin{equation}\label{eq:Precision}
\frac{1}{N} \sum_{i=0}^{N-1} \frac{\left|L_i \cap P_i\right|}{\left|P_i\right|}
\end{equation}

\textbf{Recall} is defined as
\begin{equation}\label{eq:Recall}
\frac{1}{N} \sum_{i=0}^{N-1} \frac{\left|L_i \cap P_i\right|}{\left|L_i\right|}
\end{equation}

\textbf{Accuracy} is defined as
\begin{equation}\label{eq:Accuracy}
\frac{1}{N} \sum_{i=0}^{N - 1} \frac{\left|L_i \cap P_i \right|}
        {\left|L_i\right| + \left|P_i\right| - \left|L_i \cap P_i \right|}
\end{equation}

\textbf{F1 Measure} is an integrated version combined by harmonic mean of \textbf{Precision} and \textbf{Recall}. 
\begin{equation}\label{eq:Accuracy}
\frac{1}{N} \sum_{i=0}^{N-1} 2 \frac{\left|P_i \cap L_i\right|}{\left|P_i\right| \cdot \left|L_i\right|}
\end{equation}

\subsection{Label-based Metrics}

\textbf{Macro Precision} (precision averaged across all labels) is defined as
\begin{equation}\label{eq:MacroPrecision}
PPV(\ell)=\frac{TP}{TP + FP}=
          \frac{\sum_{i=0}^{N-1} I_{P_i}(\ell) \cdot I_{L_i}(\ell)}
          {\sum_{i=0}^{N-1} I_{P_i}(\ell)}
\end{equation}

\textbf{Macro Recall} (recall averaged across all labels) is defined as
\begin{equation}\label{eq:MacroRecall}
TPR(\ell)=\frac{TP}{P}=
          \frac{\sum_{i=0}^{N-1} I_{P_i}(\ell) \cdot I_{L_i}(\ell)}
          {\sum_{i=0}^{N-1} I_{L_i}(\ell)}
\end{equation}

\textbf{F1 Measure by label} is the harmonic mean between \textbf{Precision} and \textbf{Recall}. 
\begin{equation}\label{eq:LabelAccuracy}
F1(\ell) = 2
                            \cdot \left(\frac{PPV(\ell) \cdot TPR(\ell)}
                            {PPV(\ell) + TPR(\ell)}\right)
\end{equation}

\textbf{Micro Precision} (precision averaged over all the example/label pairs) is defined as
\begin{equation}\label{eq:MicroPrecision}
\frac{TP}{TP + FP}=\frac{\sum_{i=0}^{N-1} \left|P_i \cap L_i\right|}
          {\sum_{i=0}^{N-1} \left|P_i \cap L_i\right| + \sum_{i=0}^{N-1} \left|P_i - L_i\right|}
\end{equation}

\textbf{Micro Recall} (recall averaged over all the example/label pairs) is defined as
\begin{equation}\label{eq:MicroRecall}
\frac{TP}{TP + FN}=\frac{\sum_{i=0}^{N-1} \left|P_i \cap L_i\right|}
        {\sum_{i=0}^{N-1} \left|P_i \cap L_i\right| + \sum_{i=0}^{N-1} \left|L_i - P_i\right|}
\end{equation}

\textbf{Micro F1 Measure by label} is the harmonic mean between \textbf{Micro Precision} and \textbf{Micro Recall}. 
\begin{equation}\label{eq:LabelMicroAccuracy}
2 \cdot \frac{TP}{2 \cdot TP + FP + FN}=2 \cdot \frac{\sum_{i=0}^{N-1} \left|P_i \cap L_i\right|}{2 \cdot
        \sum_{i=0}^{N-1} \left|P_i \cap L_i\right| + \sum_{i=0}^{N-1} \left|L_i - P_i\right| + \sum_{i=0}^{N-1}
        \left|P_i - L_i\right|}
\end{equation}

For the label-based metrics, the larger the metrics value means the higher generalization performance.

With the previous metrics, there are diverse methods to evaluate model generalization performance. In most multi-label training, learning algorithms optimize one of the metrics. To make evaluation fair and precise, the learning algorithms should be tested on different metrics to evaluate performance to avoid bias.

Most multi-label metrics are non-convex and discrete, most algorithms turn to optimize alternative multi-label metrics. Recent some researcher study the consistency of multi-label learning\citep{gao2013consistency}, for example, figuring if the loss of classifier converges to the Bayes loss while increasing training set size.

\section{Loss Function}

To adapt a neural network handling single-label instances to multi-label instances, we need to design a specific loss function instead of squared error loss function to track the characteristics of multi-label learning and some revision of stochastic gradient descent. 

\section{Artificial Neural Networks}



\section{Weather Classification}




